% alg-p1-01: 数轴、相反数对称、绝对值（到原点的距离）
\documentclass[tikz,border=4pt]{standalone}
\usepackage{ctex}
\usepackage{amsmath}
\usetikzlibrary{arrows.meta, decorations.pathreplacing}
\begin{document}
\begin{tikzpicture}[scale=0.85, thick]

  % 数轴主线
  \draw[-{Stealth[length=5pt]}] (-4.5,0) -- (4.8,0);

  % 刻度与标签
  \foreach \x in {-4,-3,-2,-1,0,1,2,3,4} {
    \draw (\x, 3pt) -- (\x, -3pt) node[below, font=\small] {$\x$};
  }
  \node[below right, font=\small] at (4.8,0) {};

  % 原点标 O
  \node[above, font=\small] at (0,0.08) {$O$};

  % ---- 相反数对称：-3 与 3 ----
  \filldraw[red] (-3,0) circle (2.5pt) node[above, font=\small] {$-a$};
  \filldraw[red] (3,0)  circle (2.5pt) node[above, font=\small] {$a$};

  % 对称弧（说明以 O 为中心对称）
  \draw[red, dashed, thin] (-3,0) arc (180:0:3) ;

  % 对称标注
  \node[red, font=\footnotesize] at (0, 3.3) {关于原点对称};

  % ---- 绝对值：|a| = 到原点距离 ----
  % 用花括号标出 a 到 0 的距离
  \draw[decorate,
        decoration={brace, amplitude=5pt},
        blue, thick]
    (0, -0.55) -- node[below=5pt, font=\footnotesize, blue]
    {$|a|$（到原点的距离）} (3, -0.55);

  \draw[decorate,
        decoration={brace, amplitude=5pt, mirror},
        blue, thick]
    (-3, -0.55) -- node[below=5pt, font=\footnotesize, blue]
    {$|-a| = |a|$} (0, -0.55);

\end{tikzpicture}
\end{document}
